\begin{abstract}
    Gradient-based meta-learners such as MAML adapt a shared initialization $\theta_0$ to a new task from only a handful of support examples $S_i$, but the mechanism by which $S_i$ steers this adaptation is opaque: when the support set is unrepresentative, hard, or out-of-distribution, adaptation can actively \emph{hurt} query performance, a phenomenon known as negative transfer that prior work detects only through accuracy drops, never explains at the feature level. We introduce \textbf{adaptation gain} $\Delta M$, a differentiable, sign-aware measure of how much letting the backbone adapt helps or hurts relative to a feature-reuse baseline where only the classifier head adapts, and \textbf{\fama} (Feature-space Adjoint Meta-learning Attribution), a post-hoc explainer that traces $\Delta M$ back to the support set. \fama extends Grad-CAM across an entire multi-step adaptation trajectory with the exact discrete adjoint recursion, implemented through Pearlmutter Hessian-vector products without materializing parameter-by-parameter Jacobians, and applies a signed feature-space projection at every state. We further contribute two faithfulness protocols -- a bidirectional deletion test and cascading-randomization/support-intervention sanity checks -- built for trajectories rather than static predictions. On MAML with Conv4 and ResNet-10 backbones, FAMA passes both directions of our deletion test on tieredImageNet (600 tasks, $p < 0.001$), and its maps decorrelate from the unperturbed maps once the output-nearest convolutional layer is randomized. On 1,200 cross-domain tasks (miniImageNet$\to$CUB-200-2011 and tieredImageNet$\to$CUB-200-2011), inspection of the most extreme tasks shows that negative transfer coincides with attribution shifting from the bird to background and image artifacts, a pattern consistent with \emph{contextual bias} that we present as a hypothesis rather than an established cause.
\end{abstract}
